\noindent\textbf{Coordenadas Esféricas}
\vspace{0.2cm}

Para el cálculo de integrales en regiones con simetría esférica, se utiliza el sistema de coordenadas esféricas. Un punto $(x, y, z)$ se relaciona con $(\rho, \phi, \theta)$ mediante:

\begin{equation}
    \begin{cases}
        x = \rho \sin \phi \cos \theta \\
        y = \rho \sin \phi \sin \theta \\
        z = \rho \cos \phi
    \end{cases}
\end{equation}

Donde las variables se restringen a $\rho \ge 0$, $0 \le \phi \le \pi$ y $0 \le \theta \le 2\pi$ para recorrer toda la esfera.

\vspace{0.4cm}
\noindent\textbf{El Determinante Jacobiano}

De acuerdo con el Teorema 2.5.2 de cambio de variable para integrales triples, el cambio de coordenadas cartesianas a esféricas requiere introducir el valor absoluto del determinante de la matriz Jacobiana $J(\rho, \phi, \theta)$ como factor de escala en el diferencial de volumen.

La matriz de derivadas parciales es:
\begin{equation}
J(\rho, \phi, \theta) = 
\begin{pmatrix}
\sin \phi \cos \theta & \rho \cos \phi \cos \theta & -\rho \sin \phi \sin \theta \\
\sin \phi \sin \theta & \rho \cos \phi \sin \theta & \rho \sin \phi \cos \theta \\
\cos \phi & -\rho \sin \phi & 0
\end{pmatrix}
\end{equation}

Tal como se demuestra en el citado ejemplo 2.4.5, el valor absoluto del determinante de esta matriz es:
\begin{equation}
    \left| J_{\mathbf{T}} \right| = \rho^2 \sin \phi
\end{equation}

\vspace{0.4cm}
\noindent\textbf{Diferencial de Volumen}

Dado que en el dominio de integración $\sin \phi \ge 0$, el factor de escala es simplemente $\rho^2 \sin \phi$. Por lo tanto, siguiendo el \textbf{Teorema de Cambio de Variables (2.5.2)}, el diferencial de volumen resulta:

\begin{equation}
    dV = \rho^2 \sin \phi \, d\rho \, d\phi \, d\theta
\end{equation}


\vspace{1cm}


\noindent\textbf{Coordenadas Cilíndricas}
\vspace{0.2cm}

Para el cálculo de integrales en regiones con simetría cilíndrica, se utiliza el sistema de coordenadas cilíndricas. Un punto $(x, y, z)$ se relaciona con $(r, \theta, z)$ mediante:

\begin{equation}
\begin{cases} 
x = r \cos \theta \\ 
y = r \sin \theta \\ 
z = z 
\end{cases}
\end{equation}

Donde las variables se restringen a $r \ge 0$, $0 \le \theta \le 2\pi$ y $z \in \mathbb{R}$ para recorrer todo el espacio.

\vspace{0.4cm}
\noindent\textbf{El Determinante Jacobiano}

De acuerdo con el Teorema 2.5.2 de cambio de variable para integrales triples, el cambio de coordenadas cartesianas a cilíndricas requiere introducir el valor absoluto del determinante de la matriz Jacobiana $J(r, \theta, z)$ como factor de escala en el diferencial de volumen.

La matriz de derivadas parciales es:

\begin{equation}
J(r, \theta, z) = 
\begin{pmatrix} 
\cos \theta & -r \sin \theta & 0 \\ 
\sin \theta & r \cos \theta & 0 \\ 
0 & 0 & 1 
\end{pmatrix}
\end{equation}

Tal como se demuestra en el citado ejemplo 2.4.4, el valor absoluto del determinante de esta matriz es:

\begin{equation}
|J_{\mathbf{T}}| = r
\end{equation}

\vspace{0.4cm}
\noindent\textbf{Diferencial de Volumen}

Dado que en el dominio de integración $r \ge 0$, el factor de escala es simplemente $r$. Por lo tanto, siguiendo el \textbf{Teorema de Cambio de Variables (2.5.2)}, el diferencial de volumen resulta:

\begin{equation}
dV = r \, dr \, d\theta \, dz
\end{equation}